% Generated. Do not edit.
% Projection scriptural-date-location/section, version 1: the Scriptural Date and
% Location page, generated whole from research/chronology.toml, the Scripture
% chronology corpus and the source library by tools/tpt document-types refresh.
% document: liturgy/roman-rite/postconciliar/roman-missal-third-edition-en-us-2011/propers/temporal/pc-s58-thirty-second-sunday-in-ordinary-time-year-a
\runninghead{Scriptural Date and Location}
\section{Scriptural Date and Location}
\label{sec:date-location}

{\small Each passage the Mass reads or sings from Scripture, in canonical order. The Context column names what is dated --- the event the passage narrates, a setting it remembers, the traditional attribution of its author, or the composition of its book --- and where it belongs; the Date column gives the answer of the Scripture chronology corpus for that subject, with every alternative it records. Texts of the Mass that are not drawn from Scripture have no row. The evidence for each date follows the table.\par}
\medskip
% Generated. Do not edit.
% Deterministic projection of research/chronology.toml.
% schema: 5
% document: liturgy/roman-rite/postconciliar/roman-missal-third-edition-en-us-2011/propers/temporal/pc-s58-thirty-second-sunday-in-ordinary-time-year-a
% generated-by: tools/tpt proper-chronology annotations
\newcommand{\chronologyannotationclaim}[6]{#6}
\newcommand{\chronologyannotationreach}[3]{}
\newcommand{\chronologyannotationgroup}[3]{#3}
\newcommand{\chronologyannotation}[1]{%
  \ifcsname triptychchronologyannotation@#1\endcsname
    \csname triptychchronologyannotation@#1\endcsname
  \else
    \PackageError{triptych}{No chronology annotation for #1}{}%
  \fi}
% appointment-dependency: guidance/liturgy/postconciliar-propers-registry.md sha256=376860a7c37a4c75bb6d880cd1b260183be1877a07e3cd01660a9e2854e2ea72
% appointment-dependency: src/gpt/liturgy/roman-rite/postconciliar/roman-missal-third-edition-en-us-2011/propers/registry/README.md sha256=df0aae579b2ae7e02cca5597f382ee615974bc84fc46fea909a9c21d942ce152
% appointment-dependency: src/gpt/liturgy/roman-rite/postconciliar/roman-missal-third-edition-en-us-2011/propers/registry/formula-dispositions.md sha256=af014f05723fb53a99006ef740101dc8abf629ed8b0b53a60069e94713cb289e
% appointment-dependency: src/gpt/liturgy/roman-rite/postconciliar/roman-missal-third-edition-en-us-2011/propers/temporal/pc-s58-thirty-second-sunday-in-ordinary-time-year-a/instance/manifest.md sha256=f38421227cac954eb91160f0fcbf1d8e80909e287d1beb7648bf673bfdcf6bd7
% appointment-dependency: src/gpt/liturgy/roman-rite/postconciliar/roman-missal-third-edition-en-us-2011/propers/temporal/pc-s58-thirty-second-sunday-in-ordinary-time-year-a/propers/verified.md sha256=c4a78aa4cd853e8945252cd93d625fa467fb9e340e58d1ce71fcb7e9999483a2
% appointment-dependency: src/gpt/liturgy/roman-rite/postconciliar/roman-missal-third-edition-en-us-2011/propers/temporal/pc-s58-thirty-second-sunday-in-ordinary-time-year-a/research/chronology-inputs.toml sha256=e1b26460bbd35d0bcb30648f173ce2ac5904861b663e666a1ceaf1ec5e04ee76
% appointment-dependency: src/gpt/liturgy/roman-rite/postconciliar/roman-missal-third-edition-en-us-2011/propers/temporal/shared/ordinary-time/weeks/32/propers/verified.md sha256=779f4cfc8540b75ceb3df58116c24f129230d34a33323a0e8171175cbdd36c66
% appointment-note: entrance: Cf. Psalm 87:3; adaptation; vulgate numbering; canonical verse queries, not appointed wording.
% appointment-note: collect: Composed Missal oration with no appointed biblical citation; thematic affinities create no additional Scripture appointment.
% appointment-note: first-reading: Wisdom 6:13-17; appointed; vulgate numbering; canonical verse queries, not appointed wording.
% appointment-note: psalm: Psalm 63:2,3-4,5-6,7-8; appointed; hebrew numbering; canonical verse queries, not appointed wording.
% appointment-note: psalm: Psalm 63:2b; response; hebrew numbering; whole-verse query envelope, not subverse chronology or appointed wording.
% appointment-note: second-reading-long: 1 Thessalonians 4:13-18; appointed; vulgate numbering; canonical verse queries, not appointed wording.
% appointment-note: second-reading-short: 1 Thessalonians 4:13-14; appointed; vulgate numbering; canonical verse queries, not appointed wording.
% appointment-note: acclamation: Cf. Matthew 24:42a,44; adaptation; vulgate numbering; whole-verse query envelope, not subverse chronology or appointed wording.
% appointment-note: gospel: Matthew 25:1-13; appointed; vulgate numbering; canonical verse queries, not appointed wording.
% appointment-note: offerings: Composed Missal oration with no appointed biblical citation; thematic affinities create no additional Scripture appointment.
% appointment-note: communion-psalm: Cf. Psalm 22:1-2; adaptation; vulgate numbering; canonical verse queries, not appointed wording.
% appointment-note: communion-luke: Cf. Luke 24:35; adaptation; vulgate numbering; canonical verse queries, not appointed wording.
% appointment-note: after-communion: Composed Missal oration with no appointed biblical citation; thematic affinities create no additional Scripture appointment.
% entrance: Entrance antiphon
\expandafter\def\csname triptychchronologyannotation@entrance\endcsname{%
  \chronologyannotationgroup{historical-setting}{research-pending}{\textbf{Historical setting} -- Occasion date unresolved.}%
  \space \chronologyannotationgroup{composition}{preferred}{\textbf{Composition}: \chronologyannotationreach{Ps.87.3}{Ps.41 Ps.42 Ps.43 Ps.44 Ps.45 Ps.46 Ps.47 Ps.48 Ps.83 Ps.84 Ps.86 Ps.87}{true}\chronologyannotationclaim{composition.psalms-of-the-sons-of-korah}{composition}{catholic-traditional-v1}{preferred}{between the days of Isaias and the return from exile}{Between the days of Isaias and the return from exile}.}%
}
% first-reading: First reading (appointed Wisdom 6:12--16); Vulgate equivalent
\expandafter\def\csname triptychchronologyannotation@first-reading\endcsname{%
  \chronologyannotationgroup{composition}{disputed}{\textbf{Composition} -- disputed: \chronologyannotationreach{Wis.6.13}{Wis}{true}\chronologyannotationreach{Wis.6.14}{Wis}{true}\chronologyannotationreach{Wis.6.15}{Wis}{true}\chronologyannotationreach{Wis.6.16}{Wis}{true}\chronologyannotationreach{Wis.6.17}{Wis}{true}\chronologyannotationclaim{composition.book-of-wisdom}{composition}{catholic-traditional-v1}{disputed}{the time of ... Ptolemy IV Philopator (221-204 B.C.)}{In or after the time of Ptolemy IV Philopator (221-204 B.C.)}; \chronologyannotationreach{Wis.6.13}{Wis}{true}\chronologyannotationreach{Wis.6.14}{Wis}{true}\chronologyannotationreach{Wis.6.15}{Wis}{true}\chronologyannotationreach{Wis.6.16}{Wis}{true}\chronologyannotationreach{Wis.6.17}{Wis}{true}\chronologyannotationclaim{composition.book-of-wisdom}{composition}{catholic-traditional-v1}{disputed}{the time of ... Ptolemy VII Physicon (145-117 B.C.)}{In or after the time of Ptolemy VII Physicon (145-117 B.C.)}.}%
}
% psalm: Responsorial psalm
\expandafter\def\csname triptychchronologyannotation@psalm\endcsname{%
  \chronologyannotationgroup{traditional-attribution}{disputed}{\textbf{Traditional attribution} -- disputed: \chronologyannotationreach{Ps.62.2}{Ps.22 Ps.62}{true}\chronologyannotationreach{Ps.62.3}{Ps.22 Ps.62}{true}\chronologyannotationreach{Ps.62.4}{Ps.22 Ps.62}{true}\chronologyannotationreach{Ps.62.5}{Ps.22 Ps.62}{true}\chronologyannotationreach{Ps.62.6}{Ps.22 Ps.62}{true}\chronologyannotationreach{Ps.62.7}{Ps.22 Ps.62}{true}\chronologyannotationreach{Ps.62.8}{Ps.22 Ps.62}{true}\chronologyannotationclaim{israel.monarchy.david-traditional-era}{traditional-attribution}{catholic-traditional-v1}{disputed}{reigned from 1055 to 1015 B.C.}{David (reign in the usual chronology), B.C. 1055--1015}.}%
  \space \chronologyannotationgroup{historical-setting}{preferred}{\textbf{Historical setting}: \chronologyannotationreach{Ps.62.2}{Ps.62 Ps.70}{true}\chronologyannotationreach{Ps.62.3}{Ps.62 Ps.70}{true}\chronologyannotationreach{Ps.62.4}{Ps.62 Ps.70}{true}\chronologyannotationreach{Ps.62.5}{Ps.62 Ps.70}{true}\chronologyannotationreach{Ps.62.6}{Ps.62 Ps.70}{true}\chronologyannotationreach{Ps.62.7}{Ps.62 Ps.70}{true}\chronologyannotationreach{Ps.62.8}{Ps.62 Ps.70}{true}\chronologyannotationclaim{israel.monarchy.david-flight-from-absalom}{historical-setting}{catholic-traditional-v1}{preferred}{when he fled from the face of his son Absalom}{David's flight from the face of Absalom, when he fled from the face of his son Absalom}.}%
  \space \chronologyannotationgroup{composition}{preferred}{\textbf{Composition}: \chronologyannotationreach{Ps.62.2}{Ps}{true}\chronologyannotationreach{Ps.62.3}{Ps}{true}\chronologyannotationreach{Ps.62.4}{Ps}{true}\chronologyannotationreach{Ps.62.5}{Ps}{true}\chronologyannotationreach{Ps.62.6}{Ps}{true}\chronologyannotationreach{Ps.62.7}{Ps}{true}\chronologyannotationreach{Ps.62.8}{Ps}{true}\chronologyannotationclaim{critical.psalms.latest-composition-boundary}{composition}{catholic-critical-v1}{preferred}{before the Maccabean period, around 165 B.C.; no individual psalm can be dated securely}{Before c. 165 B.C.}}%
}
% second-reading-long: Second reading: long
\expandafter\def\csname triptychchronologyannotation@second-reading-long\endcsname{%
  \chronologyannotationgroup{composition}{disputed}{\textbf{Composition} -- disputed: \chronologyannotationreach{1Thess.4.13}{1Thess}{true}\chronologyannotationreach{1Thess.4.14}{1Thess}{true}\chronologyannotationreach{1Thess.4.15}{1Thess}{true}\chronologyannotationreach{1Thess.4.16}{1Thess}{true}\chronologyannotationreach{1Thess.4.17}{1Thess}{true}\chronologyannotationreach{1Thess.4.18}{1Thess}{true}\chronologyannotationclaim{composition.first-epistle-to-the-thessalonians}{composition}{catholic-traditional-v1}{disputed}{(1 and 2 Thessalonians), 52}{A.D. 52}; \chronologyannotationreach{1Thess.4.13}{1Thess}{true}\chronologyannotationreach{1Thess.4.14}{1Thess}{true}\chronologyannotationreach{1Thess.4.15}{1Thess}{true}\chronologyannotationreach{1Thess.4.16}{1Thess}{true}\chronologyannotationreach{1Thess.4.17}{1Thess}{true}\chronologyannotationreach{1Thess.4.18}{1Thess}{true}\chronologyannotationclaim{composition.first-epistle-to-the-thessalonians}{composition}{catholic-traditional-v1}{disputed}{in the year 53 or 54 according to the commonly received scheme of Pauline chronology}{A.D. 53--54}.}%
}
% second-reading-short: Second reading: short
\expandafter\def\csname triptychchronologyannotation@second-reading-short\endcsname{%
  \chronologyannotationgroup{composition}{disputed}{\textbf{Composition} -- disputed: \chronologyannotationreach{1Thess.4.13}{1Thess}{true}\chronologyannotationreach{1Thess.4.14}{1Thess}{true}\chronologyannotationclaim{composition.first-epistle-to-the-thessalonians}{composition}{catholic-traditional-v1}{disputed}{(1 and 2 Thessalonians), 52}{A.D. 52}; \chronologyannotationreach{1Thess.4.13}{1Thess}{true}\chronologyannotationreach{1Thess.4.14}{1Thess}{true}\chronologyannotationclaim{composition.first-epistle-to-the-thessalonians}{composition}{catholic-traditional-v1}{disputed}{in the year 53 or 54 according to the commonly received scheme of Pauline chronology}{A.D. 53--54}.}%
}
% acclamation: Gospel acclamation
\expandafter\def\csname triptychchronologyannotation@acclamation\endcsname{%
  \chronologyannotationgroup{composition}{disputed}{\textbf{Composition} -- disputed: \chronologyannotationreach{Matt.24.42}{Matt}{true}\chronologyannotationreach{Matt.24.44}{Matt}{true}\chronologyannotationclaim{composition.gospel-of-matthew}{composition}{catholic-traditional-v1}{disputed}{about A.D. 38-45}{c. A.D. 38--45}; \chronologyannotationreach{Matt.24.42}{Matt}{true}\chronologyannotationreach{Matt.24.44}{Matt}{true}\chronologyannotationclaim{composition.gospel-of-matthew}{composition}{catholic-traditional-v1}{disputed}{about the year 40-42}{c. A.D. 40--42}; \chronologyannotationreach{Matt.24.42}{Matt}{true}\chronologyannotationreach{Matt.24.44}{Matt}{true}\chronologyannotationclaim{composition.gospel-of-matthew}{composition}{catholic-traditional-v1}{disputed}{the years 40-45}{A.D. 40--45}; \chronologyannotationreach{Matt.24.42}{Matt}{true}\chronologyannotationreach{Matt.24.44}{Matt}{true}\chronologyannotationclaim{composition.gospel-of-matthew}{composition}{catholic-traditional-v1}{disputed}{about the year 60-68}{c. A.D. 60--68}; \chronologyannotationreach{Matt.24.42}{Matt}{true}\chronologyannotationreach{Matt.24.44}{Matt}{true}\chronologyannotationclaim{composition.gospel-of-matthew}{composition}{catholic-traditional-v1}{disputed}{about the years 64-67}{c. A.D. 64--67}; \chronologyannotationreach{Matt.24.42}{Matt}{true}\chronologyannotationreach{Matt.24.44}{Matt}{true}\chronologyannotationclaim{composition.gospel-of-matthew}{composition}{catholic-traditional-v1}{disputed}{about the year 50}{c. A.D. 50}.}%
}
% gospel: Gospel
\expandafter\def\csname triptychchronologyannotation@gospel\endcsname{%
  \chronologyannotationgroup{narrated-event}{research-pending}{\textbf{Event} -- No narrated-event date in the chronology corpus.}%
  \space \chronologyannotationgroup{composition}{disputed}{\textbf{Composition} -- disputed: \chronologyannotationreach{Matt.25.1}{Matt}{true}\chronologyannotationreach{Matt.25.10}{Matt}{true}\chronologyannotationreach{Matt.25.11}{Matt}{true}\chronologyannotationreach{Matt.25.12}{Matt}{true}\chronologyannotationreach{Matt.25.13}{Matt}{true}\chronologyannotationreach{Matt.25.2}{Matt}{true}\chronologyannotationreach{Matt.25.3}{Matt}{true}\chronologyannotationreach{Matt.25.4}{Matt}{true}\chronologyannotationreach{Matt.25.5}{Matt}{true}\chronologyannotationreach{Matt.25.6}{Matt}{true}\chronologyannotationreach{Matt.25.7}{Matt}{true}\chronologyannotationreach{Matt.25.8}{Matt}{true}\chronologyannotationreach{Matt.25.9}{Matt}{true}\chronologyannotationclaim{composition.gospel-of-matthew}{composition}{catholic-traditional-v1}{disputed}{about A.D. 38-45}{c. A.D. 38--45}; \chronologyannotationreach{Matt.25.1}{Matt}{true}\chronologyannotationreach{Matt.25.10}{Matt}{true}\chronologyannotationreach{Matt.25.11}{Matt}{true}\chronologyannotationreach{Matt.25.12}{Matt}{true}\chronologyannotationreach{Matt.25.13}{Matt}{true}\chronologyannotationreach{Matt.25.2}{Matt}{true}\chronologyannotationreach{Matt.25.3}{Matt}{true}\chronologyannotationreach{Matt.25.4}{Matt}{true}\chronologyannotationreach{Matt.25.5}{Matt}{true}\chronologyannotationreach{Matt.25.6}{Matt}{true}\chronologyannotationreach{Matt.25.7}{Matt}{true}\chronologyannotationreach{Matt.25.8}{Matt}{true}\chronologyannotationreach{Matt.25.9}{Matt}{true}\chronologyannotationclaim{composition.gospel-of-matthew}{composition}{catholic-traditional-v1}{disputed}{about the year 40-42}{c. A.D. 40--42}; \chronologyannotationreach{Matt.25.1}{Matt}{true}\chronologyannotationreach{Matt.25.10}{Matt}{true}\chronologyannotationreach{Matt.25.11}{Matt}{true}\chronologyannotationreach{Matt.25.12}{Matt}{true}\chronologyannotationreach{Matt.25.13}{Matt}{true}\chronologyannotationreach{Matt.25.2}{Matt}{true}\chronologyannotationreach{Matt.25.3}{Matt}{true}\chronologyannotationreach{Matt.25.4}{Matt}{true}\chronologyannotationreach{Matt.25.5}{Matt}{true}\chronologyannotationreach{Matt.25.6}{Matt}{true}\chronologyannotationreach{Matt.25.7}{Matt}{true}\chronologyannotationreach{Matt.25.8}{Matt}{true}\chronologyannotationreach{Matt.25.9}{Matt}{true}\chronologyannotationclaim{composition.gospel-of-matthew}{composition}{catholic-traditional-v1}{disputed}{the years 40-45}{A.D. 40--45}; \chronologyannotationreach{Matt.25.1}{Matt}{true}\chronologyannotationreach{Matt.25.10}{Matt}{true}\chronologyannotationreach{Matt.25.11}{Matt}{true}\chronologyannotationreach{Matt.25.12}{Matt}{true}\chronologyannotationreach{Matt.25.13}{Matt}{true}\chronologyannotationreach{Matt.25.2}{Matt}{true}\chronologyannotationreach{Matt.25.3}{Matt}{true}\chronologyannotationreach{Matt.25.4}{Matt}{true}\chronologyannotationreach{Matt.25.5}{Matt}{true}\chronologyannotationreach{Matt.25.6}{Matt}{true}\chronologyannotationreach{Matt.25.7}{Matt}{true}\chronologyannotationreach{Matt.25.8}{Matt}{true}\chronologyannotationreach{Matt.25.9}{Matt}{true}\chronologyannotationclaim{composition.gospel-of-matthew}{composition}{catholic-traditional-v1}{disputed}{about the year 60-68}{c. A.D. 60--68}; \chronologyannotationreach{Matt.25.1}{Matt}{true}\chronologyannotationreach{Matt.25.10}{Matt}{true}\chronologyannotationreach{Matt.25.11}{Matt}{true}\chronologyannotationreach{Matt.25.12}{Matt}{true}\chronologyannotationreach{Matt.25.13}{Matt}{true}\chronologyannotationreach{Matt.25.2}{Matt}{true}\chronologyannotationreach{Matt.25.3}{Matt}{true}\chronologyannotationreach{Matt.25.4}{Matt}{true}\chronologyannotationreach{Matt.25.5}{Matt}{true}\chronologyannotationreach{Matt.25.6}{Matt}{true}\chronologyannotationreach{Matt.25.7}{Matt}{true}\chronologyannotationreach{Matt.25.8}{Matt}{true}\chronologyannotationreach{Matt.25.9}{Matt}{true}\chronologyannotationclaim{composition.gospel-of-matthew}{composition}{catholic-traditional-v1}{disputed}{about the years 64-67}{c. A.D. 64--67}; \chronologyannotationreach{Matt.25.1}{Matt}{true}\chronologyannotationreach{Matt.25.10}{Matt}{true}\chronologyannotationreach{Matt.25.11}{Matt}{true}\chronologyannotationreach{Matt.25.12}{Matt}{true}\chronologyannotationreach{Matt.25.13}{Matt}{true}\chronologyannotationreach{Matt.25.2}{Matt}{true}\chronologyannotationreach{Matt.25.3}{Matt}{true}\chronologyannotationreach{Matt.25.4}{Matt}{true}\chronologyannotationreach{Matt.25.5}{Matt}{true}\chronologyannotationreach{Matt.25.6}{Matt}{true}\chronologyannotationreach{Matt.25.7}{Matt}{true}\chronologyannotationreach{Matt.25.8}{Matt}{true}\chronologyannotationreach{Matt.25.9}{Matt}{true}\chronologyannotationclaim{composition.gospel-of-matthew}{composition}{catholic-traditional-v1}{disputed}{about the year 50}{c. A.D. 50}.}%
}
% communion-psalm: Communion: Psalm
\expandafter\def\csname triptychchronologyannotation@communion-psalm\endcsname{%
  \chronologyannotationgroup{traditional-attribution}{disputed}{\textbf{Traditional attribution} -- disputed: \chronologyannotationreach{Ps.22.1}{Ps.22 Ps.62}{true}\chronologyannotationreach{Ps.22.2}{Ps.22 Ps.62}{true}\chronologyannotationclaim{israel.monarchy.david-traditional-era}{traditional-attribution}{catholic-traditional-v1}{disputed}{reigned from 1055 to 1015 B.C.}{David (reign in the usual chronology), B.C. 1055--1015}.}%
  \space \chronologyannotationgroup{historical-setting}{research-pending}{\textbf{Historical setting} -- Occasion date unresolved.}%
  \space \chronologyannotationgroup{composition}{preferred}{\textbf{Composition}: \chronologyannotationreach{Ps.22.1}{Ps}{true}\chronologyannotationreach{Ps.22.2}{Ps}{true}\chronologyannotationclaim{critical.psalms.latest-composition-boundary}{composition}{catholic-critical-v1}{preferred}{before the Maccabean period, around 165 B.C.; no individual psalm can be dated securely}{Before c. 165 B.C.}}%
}
% communion-luke: Communion: Luke
\expandafter\def\csname triptychchronologyannotation@communion-luke\endcsname{%
  \chronologyannotationgroup{composition}{disputed}{\textbf{Composition} -- disputed: \chronologyannotationreach{Luke.24.35}{Luke}{true}\chronologyannotationclaim{composition.gospel-of-luke}{composition}{catholic-traditional-v1}{disputed}{About the year 70}{c. A.D. 70}; \chronologyannotationreach{Luke.24.35}{Luke}{true}\chronologyannotationclaim{composition.gospel-of-luke}{composition}{catholic-traditional-v1}{disputed}{before the end of the Roman imprisonment, when the Acts was finished}{Before the end of the Roman imprisonment, when the Acts was finished}.}%
}
% presentation: orientation-v3
\newcommand{\chronologyorientationrow}[2]{%
  \ifcsname triptychchronologyorientation@#1\endcsname
    \csname triptychchronologyorientation@#1\endcsname{#2}%
  \else
    \PackageError{triptych}{No chronology orientation row for #1}{}%
  \fi}
\expandafter\def\csname triptychchronologyorientation@row-85017d44138f4354\endcsname#1{%
Communion: Psalm: Cf. Ps. 22:1--2 (Vg.; Heb. 23) & Traditional attribution: David's reign in the usual chronology; Hebron, then Jerusalem & Disputed: \chronologyannotationreach{Ps.22.1}{Ps.22 Ps.62}{true}\chronologyannotationreach{Ps.22.2}{Ps.22 Ps.62}{true}\chronologyannotationclaim{israel.monarchy.david-traditional-era}{traditional-attribution}{catholic-traditional-v1}{disputed}{reigned from 1055 to 1015 B.C.}{B.C.~1055--1015 (Corbett also reports writers dating the reign 30--50 years later)}\\*[.1em]%
}
\expandafter\def\csname triptychchronologyorientation@row-4f6da3e110136fb3\endcsname#1{%
 & Historical setting & Occasion date unresolved\\*[.1em]%
}
\expandafter\def\csname triptychchronologyorientation@row-53baae847fde63b6\endcsname#1{%
 & Composition: Psalms & \chronologyannotationreach{Ps.22.1}{Ps}{true}\chronologyannotationreach{Ps.22.2}{Ps}{true}\chronologyannotationclaim{critical.psalms.latest-composition-boundary}{composition}{catholic-critical-v1}{preferred}{before the Maccabean period, around 165 B.C.; no individual psalm can be dated securely}{Before c.~165 B.C.}\\[.3em]%
}
\expandafter\def\csname triptychchronologyorientation@row-b79aa84c431dec69\endcsname#1{%
Responsorial psalm: Ps. 63:2,3--4,5--6,7--8 (Heb.; Vg. 62), Ps. 63:2b (Heb.; Vg. 62) & Traditional attribution: David's reign in the usual chronology; Hebron, then Jerusalem & Disputed: \chronologyannotationreach{Ps.62.2}{Ps.22 Ps.62}{true}\chronologyannotationreach{Ps.62.3}{Ps.22 Ps.62}{true}\chronologyannotationreach{Ps.62.4}{Ps.22 Ps.62}{true}\chronologyannotationreach{Ps.62.5}{Ps.22 Ps.62}{true}\chronologyannotationreach{Ps.62.6}{Ps.22 Ps.62}{true}\chronologyannotationreach{Ps.62.7}{Ps.22 Ps.62}{true}\chronologyannotationreach{Ps.62.8}{Ps.22 Ps.62}{true}\chronologyannotationclaim{israel.monarchy.david-traditional-era}{traditional-attribution}{catholic-traditional-v1}{disputed}{reigned from 1055 to 1015 B.C.}{B.C.~1055--1015 (Corbett also reports writers dating the reign 30--50 years later)}\\*[.1em]%
}
\expandafter\def\csname triptychchronologyorientation@row-eb1055c5f42de2df\endcsname#1{%
 & Historical setting: David's flight from the face of Absalom & \chronologyannotationreach{Ps.62.2}{Ps.62 Ps.70}{true}\chronologyannotationreach{Ps.62.3}{Ps.62 Ps.70}{true}\chronologyannotationreach{Ps.62.4}{Ps.62 Ps.70}{true}\chronologyannotationreach{Ps.62.5}{Ps.62 Ps.70}{true}\chronologyannotationreach{Ps.62.6}{Ps.62 Ps.70}{true}\chronologyannotationreach{Ps.62.7}{Ps.62 Ps.70}{true}\chronologyannotationreach{Ps.62.8}{Ps.62 Ps.70}{true}\chronologyannotationclaim{israel.monarchy.david-flight-from-absalom}{historical-setting}{catholic-traditional-v1}{preferred}{when he fled from the face of his son Absalom}{When he fled from the face of his son Absalom}\\*[.1em]%
}
\expandafter\def\csname triptychchronologyorientation@row-241ec7af94bd3fd9\endcsname#1{%
 & Composition: Psalms & \chronologyannotationreach{Ps.62.2}{Ps}{true}\chronologyannotationreach{Ps.62.3}{Ps}{true}\chronologyannotationreach{Ps.62.4}{Ps}{true}\chronologyannotationreach{Ps.62.5}{Ps}{true}\chronologyannotationreach{Ps.62.6}{Ps}{true}\chronologyannotationreach{Ps.62.7}{Ps}{true}\chronologyannotationreach{Ps.62.8}{Ps}{true}\chronologyannotationclaim{critical.psalms.latest-composition-boundary}{composition}{catholic-critical-v1}{preferred}{before the Maccabean period, around 165 B.C.; no individual psalm can be dated securely}{Before c.~165 B.C.}\\[.3em]%
}
\expandafter\def\csname triptychchronologyorientation@row-60bba31394be5ffc\endcsname#1{%
Entrance antiphon: Cf. Ps. 87:3 (Vg.; Heb. 88) & Historical setting & Occasion date unresolved\\*[.1em]%
}
\expandafter\def\csname triptychchronologyorientation@row-bddc3f22bf76bc58\endcsname#1{%
 & Composition: The Psalms of the Sons of Korah & \chronologyannotationreach{Ps.87.3}{Ps.41 Ps.42 Ps.43 Ps.44 Ps.45 Ps.46 Ps.47 Ps.48 Ps.83 Ps.84 Ps.86 Ps.87}{true}\chronologyannotationclaim{composition.psalms-of-the-sons-of-korah}{composition}{catholic-traditional-v1}{preferred}{between the days of Isaias and the return from exile}{Between the days of Isaias and the return from exile}\\[.3em]%
}
\expandafter\def\csname triptychchronologyorientation@row-34bb557c92c501a0\endcsname#1{%
First reading (appointed Wisdom 6:12--16); Vulgate equivalent: Wis. 6:13--17 & Composition: Wisdom; Alexandria in Gigot's account & Disputed: \chronologyannotationreach{Wis.6.13}{Wis}{true}\chronologyannotationreach{Wis.6.14}{Wis}{true}\chronologyannotationreach{Wis.6.15}{Wis}{true}\chronologyannotationreach{Wis.6.16}{Wis}{true}\chronologyannotationreach{Wis.6.17}{Wis}{true}\chronologyannotationclaim{composition.book-of-wisdom}{composition}{catholic-traditional-v1}{disputed}{the time of ... Ptolemy IV Philopator (221-204 B.C.)}{In or after the time of Ptolemy IV Philopator (221--204 B.C.)}; \chronologyannotationreach{Wis.6.13}{Wis}{true}\chronologyannotationreach{Wis.6.14}{Wis}{true}\chronologyannotationreach{Wis.6.15}{Wis}{true}\chronologyannotationreach{Wis.6.16}{Wis}{true}\chronologyannotationreach{Wis.6.17}{Wis}{true}\chronologyannotationclaim{composition.book-of-wisdom}{composition}{catholic-traditional-v1}{disputed}{the time of ... Ptolemy VII Physicon (145-117 B.C.)}{In or after the time of Ptolemy VII Physicon (145--117 B.C.)}\\[.3em]%
}
\expandafter\def\csname triptychchronologyorientation@row-9680220f312b5308\endcsname#1{%
Gospel acclamation: Cf. Matt. 24:42a,44 & Composition: Matthew; Semitic original associated with Palestine & Disputed: \chronologyannotationreach{Matt.24.42}{Matt}{true}\chronologyannotationreach{Matt.24.44}{Matt}{true}\chronologyannotationclaim{composition.gospel-of-matthew}{composition}{catholic-traditional-v1}{disputed}{about A.D. 38-45}{c.~A.D.~38--45}; \chronologyannotationreach{Matt.24.42}{Matt}{true}\chronologyannotationreach{Matt.24.44}{Matt}{true}\chronologyannotationclaim{composition.gospel-of-matthew}{composition}{catholic-traditional-v1}{disputed}{about the year 40-42}{c.~A.D.~40--42}; \chronologyannotationreach{Matt.24.42}{Matt}{true}\chronologyannotationreach{Matt.24.44}{Matt}{true}\chronologyannotationclaim{composition.gospel-of-matthew}{composition}{catholic-traditional-v1}{disputed}{the years 40-45}{A.D.~40--45}; \chronologyannotationreach{Matt.24.42}{Matt}{true}\chronologyannotationreach{Matt.24.44}{Matt}{true}\chronologyannotationclaim{composition.gospel-of-matthew}{composition}{catholic-traditional-v1}{disputed}{about the year 60-68}{c.~A.D.~60--68}; \chronologyannotationreach{Matt.24.42}{Matt}{true}\chronologyannotationreach{Matt.24.44}{Matt}{true}\chronologyannotationclaim{composition.gospel-of-matthew}{composition}{catholic-traditional-v1}{disputed}{about the years 64-67}{c.~A.D.~64--67}; \chronologyannotationreach{Matt.24.42}{Matt}{true}\chronologyannotationreach{Matt.24.44}{Matt}{true}\chronologyannotationclaim{composition.gospel-of-matthew}{composition}{catholic-traditional-v1}{disputed}{about the year 50}{c.~A.D.~50}\\[.3em]%
}
\expandafter\def\csname triptychchronologyorientation@row-9a58e0defb61e161\endcsname#1{%
Gospel: Matt. 25:1--13 & Narrated event: Jesus tells the ten virgins parable; Mount Olivet by narrative continuity & Date unresolved\\*[.1em]%
}
\expandafter\def\csname triptychchronologyorientation@row-350a2728de63efa5\endcsname#1{%
 & Composition: Matthew; Semitic original associated with Palestine & Disputed: \chronologyannotationreach{Matt.25.1}{Matt}{true}\chronologyannotationreach{Matt.25.10}{Matt}{true}\chronologyannotationreach{Matt.25.11}{Matt}{true}\chronologyannotationreach{Matt.25.12}{Matt}{true}\chronologyannotationreach{Matt.25.13}{Matt}{true}\chronologyannotationreach{Matt.25.2}{Matt}{true}\chronologyannotationreach{Matt.25.3}{Matt}{true}\chronologyannotationreach{Matt.25.4}{Matt}{true}\chronologyannotationreach{Matt.25.5}{Matt}{true}\chronologyannotationreach{Matt.25.6}{Matt}{true}\chronologyannotationreach{Matt.25.7}{Matt}{true}\chronologyannotationreach{Matt.25.8}{Matt}{true}\chronologyannotationreach{Matt.25.9}{Matt}{true}\chronologyannotationclaim{composition.gospel-of-matthew}{composition}{catholic-traditional-v1}{disputed}{about A.D. 38-45}{c.~A.D.~38--45}; \chronologyannotationreach{Matt.25.1}{Matt}{true}\chronologyannotationreach{Matt.25.10}{Matt}{true}\chronologyannotationreach{Matt.25.11}{Matt}{true}\chronologyannotationreach{Matt.25.12}{Matt}{true}\chronologyannotationreach{Matt.25.13}{Matt}{true}\chronologyannotationreach{Matt.25.2}{Matt}{true}\chronologyannotationreach{Matt.25.3}{Matt}{true}\chronologyannotationreach{Matt.25.4}{Matt}{true}\chronologyannotationreach{Matt.25.5}{Matt}{true}\chronologyannotationreach{Matt.25.6}{Matt}{true}\chronologyannotationreach{Matt.25.7}{Matt}{true}\chronologyannotationreach{Matt.25.8}{Matt}{true}\chronologyannotationreach{Matt.25.9}{Matt}{true}\chronologyannotationclaim{composition.gospel-of-matthew}{composition}{catholic-traditional-v1}{disputed}{about the year 40-42}{c.~A.D.~40--42}; \chronologyannotationreach{Matt.25.1}{Matt}{true}\chronologyannotationreach{Matt.25.10}{Matt}{true}\chronologyannotationreach{Matt.25.11}{Matt}{true}\chronologyannotationreach{Matt.25.12}{Matt}{true}\chronologyannotationreach{Matt.25.13}{Matt}{true}\chronologyannotationreach{Matt.25.2}{Matt}{true}\chronologyannotationreach{Matt.25.3}{Matt}{true}\chronologyannotationreach{Matt.25.4}{Matt}{true}\chronologyannotationreach{Matt.25.5}{Matt}{true}\chronologyannotationreach{Matt.25.6}{Matt}{true}\chronologyannotationreach{Matt.25.7}{Matt}{true}\chronologyannotationreach{Matt.25.8}{Matt}{true}\chronologyannotationreach{Matt.25.9}{Matt}{true}\chronologyannotationclaim{composition.gospel-of-matthew}{composition}{catholic-traditional-v1}{disputed}{the years 40-45}{A.D.~40--45}; \chronologyannotationreach{Matt.25.1}{Matt}{true}\chronologyannotationreach{Matt.25.10}{Matt}{true}\chronologyannotationreach{Matt.25.11}{Matt}{true}\chronologyannotationreach{Matt.25.12}{Matt}{true}\chronologyannotationreach{Matt.25.13}{Matt}{true}\chronologyannotationreach{Matt.25.2}{Matt}{true}\chronologyannotationreach{Matt.25.3}{Matt}{true}\chronologyannotationreach{Matt.25.4}{Matt}{true}\chronologyannotationreach{Matt.25.5}{Matt}{true}\chronologyannotationreach{Matt.25.6}{Matt}{true}\chronologyannotationreach{Matt.25.7}{Matt}{true}\chronologyannotationreach{Matt.25.8}{Matt}{true}\chronologyannotationreach{Matt.25.9}{Matt}{true}\chronologyannotationclaim{composition.gospel-of-matthew}{composition}{catholic-traditional-v1}{disputed}{about the year 60-68}{c.~A.D.~60--68}; \chronologyannotationreach{Matt.25.1}{Matt}{true}\chronologyannotationreach{Matt.25.10}{Matt}{true}\chronologyannotationreach{Matt.25.11}{Matt}{true}\chronologyannotationreach{Matt.25.12}{Matt}{true}\chronologyannotationreach{Matt.25.13}{Matt}{true}\chronologyannotationreach{Matt.25.2}{Matt}{true}\chronologyannotationreach{Matt.25.3}{Matt}{true}\chronologyannotationreach{Matt.25.4}{Matt}{true}\chronologyannotationreach{Matt.25.5}{Matt}{true}\chronologyannotationreach{Matt.25.6}{Matt}{true}\chronologyannotationreach{Matt.25.7}{Matt}{true}\chronologyannotationreach{Matt.25.8}{Matt}{true}\chronologyannotationreach{Matt.25.9}{Matt}{true}\chronologyannotationclaim{composition.gospel-of-matthew}{composition}{catholic-traditional-v1}{disputed}{about the years 64-67}{c.~A.D.~64--67}; \chronologyannotationreach{Matt.25.1}{Matt}{true}\chronologyannotationreach{Matt.25.10}{Matt}{true}\chronologyannotationreach{Matt.25.11}{Matt}{true}\chronologyannotationreach{Matt.25.12}{Matt}{true}\chronologyannotationreach{Matt.25.13}{Matt}{true}\chronologyannotationreach{Matt.25.2}{Matt}{true}\chronologyannotationreach{Matt.25.3}{Matt}{true}\chronologyannotationreach{Matt.25.4}{Matt}{true}\chronologyannotationreach{Matt.25.5}{Matt}{true}\chronologyannotationreach{Matt.25.6}{Matt}{true}\chronologyannotationreach{Matt.25.7}{Matt}{true}\chronologyannotationreach{Matt.25.8}{Matt}{true}\chronologyannotationreach{Matt.25.9}{Matt}{true}\chronologyannotationclaim{composition.gospel-of-matthew}{composition}{catholic-traditional-v1}{disputed}{about the year 50}{c.~A.D.~50}\\[.3em]%
}
\expandafter\def\csname triptychchronologyorientation@row-3b13838d29e99b34\endcsname#1{%
Communion: Luke: Cf. Luke 24:35 & Composition: The Gospel of St Luke & Disputed: \chronologyannotationreach{Luke.24.35}{Luke}{true}\chronologyannotationclaim{composition.gospel-of-luke}{composition}{catholic-traditional-v1}{disputed}{About the year 70}{c.~A.D.~70}; \chronologyannotationreach{Luke.24.35}{Luke}{true}\chronologyannotationclaim{composition.gospel-of-luke}{composition}{catholic-traditional-v1}{disputed}{before the end of the Roman imprisonment, when the Acts was finished}{Before the end of the Roman imprisonment, when the Acts was finished}\\[.3em]%
}
\expandafter\def\csname triptychchronologyorientation@row-0d44e9c16cb25df7\endcsname#1{%
Second reading: long: 1 Thess. 4:13--18 & Composition: First Thessalonians; Corinth to Thessalonica & Disputed: \chronologyannotationreach{1Thess.4.13}{1Thess}{true}\chronologyannotationreach{1Thess.4.14}{1Thess}{true}\chronologyannotationreach{1Thess.4.15}{1Thess}{true}\chronologyannotationreach{1Thess.4.16}{1Thess}{true}\chronologyannotationreach{1Thess.4.17}{1Thess}{true}\chronologyannotationreach{1Thess.4.18}{1Thess}{true}\chronologyannotationclaim{composition.first-epistle-to-the-thessalonians}{composition}{catholic-traditional-v1}{disputed}{(1 and 2 Thessalonians), 52}{A.D.~52}; \chronologyannotationreach{1Thess.4.13}{1Thess}{true}\chronologyannotationreach{1Thess.4.14}{1Thess}{true}\chronologyannotationreach{1Thess.4.15}{1Thess}{true}\chronologyannotationreach{1Thess.4.16}{1Thess}{true}\chronologyannotationreach{1Thess.4.17}{1Thess}{true}\chronologyannotationreach{1Thess.4.18}{1Thess}{true}\chronologyannotationclaim{composition.first-epistle-to-the-thessalonians}{composition}{catholic-traditional-v1}{disputed}{in the year 53 or 54 according to the commonly received scheme of Pauline chronology}{A.D.~53--54}\\[.3em]%
}
\expandafter\def\csname triptychchronologyorientation@row-7f7959c8dd983b9e\endcsname#1{%
Second reading: short: 1 Thess. 4:13--14 & Composition: First Thessalonians; Corinth to Thessalonica & Disputed: \chronologyannotationreach{1Thess.4.13}{1Thess}{true}\chronologyannotationreach{1Thess.4.14}{1Thess}{true}\chronologyannotationclaim{composition.first-epistle-to-the-thessalonians}{composition}{catholic-traditional-v1}{disputed}{(1 and 2 Thessalonians), 52}{A.D.~52}; \chronologyannotationreach{1Thess.4.13}{1Thess}{true}\chronologyannotationreach{1Thess.4.14}{1Thess}{true}\chronologyannotationclaim{composition.first-epistle-to-the-thessalonians}{composition}{catholic-traditional-v1}{disputed}{in the year 53 or 54 according to the commonly received scheme of Pauline chronology}{A.D.~53--54}\\[.3em]%
}

% Generated. Do not edit.
% Complete orientation-v3 table; wording, dates and order belong to the corpus.
\begin{chronologysharedtable}
\chronologyorientationrow{row-85017d44138f4354}{}
\chronologyorientationrow{row-4f6da3e110136fb3}{}
\chronologyorientationrow{row-53baae847fde63b6}{}
\midrule
\chronologyorientationrow{row-b79aa84c431dec69}{}
\chronologyorientationrow{row-eb1055c5f42de2df}{}
\chronologyorientationrow{row-241ec7af94bd3fd9}{}
\midrule
\chronologyorientationrow{row-60bba31394be5ffc}{}
\chronologyorientationrow{row-bddc3f22bf76bc58}{}
\midrule
\chronologyorientationrow{row-34bb557c92c501a0}{}
\midrule
\chronologyorientationrow{row-9680220f312b5308}{}
\midrule
\chronologyorientationrow{row-9a58e0defb61e161}{}
\chronologyorientationrow{row-350a2728de63efa5}{}
\midrule
\chronologyorientationrow{row-3b13838d29e99b34}{}
\midrule
\chronologyorientationrow{row-0d44e9c16cb25df7}{}
\midrule
\chronologyorientationrow{row-7f7959c8dd983b9e}{}
\end{chronologysharedtable}
{\footnotesize Psalm numbering: Vg. = Vulgate; Heb. = Hebrew.\par}


\subsection{Evidence for each date}
\label{sec:date-location-evidence}
{\small For each passage, as the Scriptural Date and Location table names it: what is dated, how the corpus weighs each answer, the source's own words, and the sources cited. A subject the corpus does not date is reported as such and given no date.\par}
{\small
\par\smallskip\noindent\textit{Communion: Psalm: Cf. Ps. 22:1--2 (Vg.; Heb. 23)}\par\nobreak
\begin{itemize}[leftmargin=1.2em,itemsep=1pt,topsep=1pt,parsep=0pt]
\item \textbf{Traditional attribution} (disputed). David's reign in the usual chronology: B.C.~1055--1015 (Corbett also reports writers dating the reign 30--50 years later) (``reigned from 1055 to 1015 B.C.''; The Catholic Encyclopedia, Volume 4 (New York, 1908), article 'David, King,' printed p. 642, opening paragraph). Place: Hebron, then Jerusalem (The Holy Bible, Douay-Rheims (Challoner revision), Project Gutenberg eBook 1581: 2 Kings 5:5).
\item \textbf{Historical setting}. No date: the chronology corpus has not yet dated this (research pending).
\item \textbf{Composition} (preferred). Psalms: Before c.~165 B.C. (``before the Maccabean period, around 165 B.C.; no individual psalm can be dated securely''; Catholic critical chronology; New American Bible, Revised Edition --- USCCB web state (2026), Psalms introduction).
\end{itemize}
\par\smallskip\noindent\textit{Responsorial psalm: Ps. 63:2,3--4,5--6,7--8 (Heb.; Vg. 62), Ps. 63:2b (Heb.; Vg. 62)}\par\nobreak
\begin{itemize}[leftmargin=1.2em,itemsep=1pt,topsep=1pt,parsep=0pt]
\item \textbf{Traditional attribution}. As under Communion: Psalm: Cf. Ps. 22:1--2 (Vg.; Heb. 23).
\item \textbf{Historical setting} (preferred). David's flight from the face of Absalom: When he fled from the face of his son Absalom (The Holy Bible, Douay-Rheims (Challoner revision), Project Gutenberg eBook 1581: Ps. 3:1).
\item \textbf{Composition}. As under Communion: Psalm: Cf. Ps. 22:1--2 (Vg.; Heb. 23).
\end{itemize}
\par\smallskip\noindent\textit{Entrance antiphon: Cf. Ps. 87:3 (Vg.; Heb. 88)}\par\nobreak
\begin{itemize}[leftmargin=1.2em,itemsep=1pt,topsep=1pt,parsep=0pt]
\item \textbf{Historical setting}. As under Communion: Psalm: Cf. Ps. 22:1--2 (Vg.; Heb. 23).
\item \textbf{Composition} (preferred). The Psalms of the Sons of Korah: Between the days of Isaias and the return from exile (The Catholic Encyclopedia, Volume 12 (1911)).
\end{itemize}
\par\smallskip\noindent\textit{First reading (appointed Wisdom 6:12--16); Vulgate equivalent: Wis. 6:13--17}\par\nobreak
\begin{itemize}[leftmargin=1.2em,itemsep=1pt,topsep=1pt,parsep=0pt]
\item \textbf{Composition} (disputed). The Book of Wisdom: In or after the time of Ptolemy IV Philopator (221--204 B.C.) (``the time of ... Ptolemy IV Philopator (221-204 B.C.)''; The Catholic Encyclopedia, Volume 15 (1912)); In or after the time of Ptolemy VII Physicon (145--117 B.C.) (``the time of ... Ptolemy VII Physicon (145-117 B.C.)''; The Catholic Encyclopedia, Volume 15 (1912)). Place: Alexandria in Gigot's account (The Catholic Encyclopedia, Volume 15 (1912)).
\end{itemize}
\par\smallskip\noindent\textit{Gospel acclamation: Cf. Matt. 24:42a,44}\par\nobreak
\begin{itemize}[leftmargin=1.2em,itemsep=1pt,topsep=1pt,parsep=0pt]
\item \textbf{Composition} (disputed). The Gospel of St Matthew: c.~A.D.~38--45 (``about A.D. 38-45''; The Catholic Encyclopedia, Volume 10 (1911); The Catholic Encyclopedia, Volume 14 (New York, 1912)); c.~A.D.~40--42 (``about the year 40-42''; The Catholic Encyclopedia, Volume 10 (1911)); A.D.~40--45 (``the years 40-45''; The Catholic Encyclopedia, Volume 10 (1911)); c.~A.D.~60--68 (``about the year 60-68''; The Catholic Encyclopedia, Volume 10 (1911)); c.~A.D.~64--67 (``about the years 64-67''; The Catholic Encyclopedia, Volume 10 (1911)); c.~A.D.~50 (``about the year 50''; The Catholic Encyclopedia, Volume 14 (New York, 1912)). Place: Semitic original associated with Palestine (The Catholic Encyclopedia, Volume 10 (1911)).
\end{itemize}
\par\smallskip\noindent\textit{Gospel: Matt. 25:1--13}\par\nobreak
\begin{itemize}[leftmargin=1.2em,itemsep=1pt,topsep=1pt,parsep=0pt]
\item \textbf{Narrated event}. Jesus tells the ten virgins parable; Mount Olivet by narrative continuity (The Holy Bible, Douay-Rheims (Challoner revision), Project Gutenberg eBook 1581: Matt. 24:3, 25:1, 25:13). No date: the chronology corpus has not yet dated this (research pending).
\item \textbf{Composition}. As under Gospel acclamation: Cf. Matt. 24:42a,44.
\end{itemize}
\par\smallskip\noindent\textit{Communion: Luke: Cf. Luke 24:35}\par\nobreak
\begin{itemize}[leftmargin=1.2em,itemsep=1pt,topsep=1pt,parsep=0pt]
\item \textbf{Composition} (disputed). The Gospel of St Luke: c.~A.D.~70 (``About the year 70''; The Catholic Encyclopedia, Volume 14 (New York, 1912)); Before the end of the Roman imprisonment, when the Acts was finished (De auctore, de tempore compositionis et de historica veritate Evangeliorum secundum Marcum et secundum Lucam (1912), responsa-i-ix).
\end{itemize}
\par\smallskip\noindent\textit{Second reading: long: 1 Thess. 4:13--18}\par\nobreak
\begin{itemize}[leftmargin=1.2em,itemsep=1pt,topsep=1pt,parsep=0pt]
\item \textbf{Composition} (disputed). The First Epistle to the Thessalonians: A.D.~52 (``(1 and 2 Thessalonians), 52''; The Catholic Encyclopedia, Volume 11 (New York, 1911)); A.D.~53--54 (``in the year 53 or 54 according to the commonly received scheme of Pauline chronology''; The Catholic Encyclopedia, Volume 14 (New York, 1912)). Place: Corinth to Thessalonica (The Catholic Encyclopedia, Volume 14 (New York, 1912); The Holy Bible, Douay-Rheims (Challoner revision), Project Gutenberg eBook 1581: 1 Thess. 1:1).
\end{itemize}
\par\smallskip\noindent\textit{Second reading: short: 1 Thess. 4:13--14}\par\nobreak
\begin{itemize}[leftmargin=1.2em,itemsep=1pt,topsep=1pt,parsep=0pt]
\item \textbf{Composition}. As under Second reading: long: 1 Thess. 4:13--18.
\end{itemize}
\par}
